\section*{Devoir maison }

On s'intéresse au code suivant qui représente un paquet de carte de jeu classique.
Un jeu classique de 52 cartes est organisé en quatre couleurs : trèfle (noir), carreau (rouge), cœur (rouge), pique (noir) réparties en quatre familles de treize (As, 2, 3, 4, 5, 6, 7, 8, 9, 10, Valet, Dame, Roi).\\

Les enseignes sont plus généralement appelées couleurs. 
\textbf{Dans le contexte du jeu de cartes, le terme couleur n'a donc pas le même sens que dans le langage commun} (dans ce dernier, on dirait qu'il n'y a que deux couleurs : le rouge et le noir).
\vfill
\begin{lstlisting}[language=Python]
class Carte:

    def __init__(self, c, v):
    """Initialise couleur (entre 1 a 4), et valeur (entre 1 a 13)"""
        self.couleur = c
        self.valeur = v

    def get_nom(self):
    """Renvoie le nom de la Carte As, 2, ... 10, Valet, Dame, Roi"""
        if self.valeur > 1 and self.valeur < 11:
            return str( self.valeur)
        elif self.valeur == 11:
            return "Valet"
        elif self.valeur == 12:
            return "Dame"
        elif self.valeur == 13:
            return "Roi"
        else:
            return "As"

    def get_couleur(self):
    """Renvoie la couleur (parmi pique, coeur, carreau, trefle"""
        return ['pique', 'coeur', 'carreau', 'trefle' ][self.couleur - 1]

class Paquet:

    def __init__(self):
        self.contenu = []

    def remplir(self):
        """Remplit le paquet de cartes"""
        for couleur in range(1, ...):
            for valeur in range( 1, ...):
                ..........................

    def get_carte_at(self, pos):
        """Renvoie la Carte qui se trouve a  la position donnee"""
        if 0 <= pos < .......................................... :
            return ...................

    def nb_cartes(self):
        return len(self.contenu)
\end{lstlisting}
\begin{Exercise}[title={POO - Cartes}]
    \Question Quels sont les attributs de la classe Carte ?
    \Question Quelles sont les méthodes de la classe Carte ?
    \begin{lstlisting}[language=Python]
carte1 = Carte(2, 2)
print(carte1.valeur)
print(carte1.couleur)
print(carte1.get_nom())
print(caret1.get_couleur())
carte2 = Carte(12, 3)
print(carte2.valeur)
print(carte2.couleur)
print(carte2.get_nom())
print(caret2.get_couleur())
    \end{lstlisting}
    \Question Que va afficher le code ci-dessus ?\\
    \emptybox{4cm}
    \Question Quels sont les attributs de la classe Paquet ?
    \Question Quelles sont les méthodes de la classe Paquet ?
    \Question Compléter la méthode remplir.
    \Question Compléter la méthode get\_carte\_at.
    \begin{lstlisting}[language=Python]
carte1 = Carte(2, 2)
paquet = Paquet()
paquet.remplir()
for i in range(1, len(paquet.nb_cartes)+1):
    current = paquet.get_carte_at(i)
    if  current.couleur == carte1.couleur and current.valeur == carte1.valeur :
        print("Trouvé en position {i}")
    \end{lstlisting}
    \Question Que va afficher le code ci-dessus ?\\
    \emptybox{1cm}
    \Question Pourquoi ne peut-on pas directement comparer current et carte1 ?\\
    \emptybox{3cm}
\end{Exercise}